
%\needspace{4\baselineskip}
\color{uniblue}{\subsection{Устройство отключения нагрузки (УОН)}}
\color{black}

\begin{enumerate}
%%%%%%%%%%%%%%%%%%%%%%%%%%%%%%%%%%%%%%%%%%%%%%%%%%%%%%%%%%%  ОБЩИЕ СВЕДЕНИЯ %%%%%%%%%%%%%%%%%%%%%%%%%%%%%%%%%%%%%%%%%%%%%%%%

\item Общие сведения

\begin{enumerate}

\item
УОН предназначено для выдачи команд на отключение или включение нагрузки потребителей от внешнего сигнала, как правило, для реализации управляющего воздействия от ПА с целью предотвращения нарушения устойчивости и разделения энергосистемы на несинхронные части.

\item 
В устройстве УОН представлено функциональным блоком «УОН» .Функциональный блок включает в себя следующие функции:
\begin{itemize}
    \item логика реализации управляющих воздействий на отключение нагрузки (УВ ОН1 ... УВ~ОН10);
    \item логика реализации управляющих воздействий на централизованное включение нагрузки (УВ ЦВН1 ... УВ~ЦВН10);
    \item сборка управляющих воздействий УОН.   
\end{itemize}

\item
УОН позволяет коммутировать входящие команды отключения и включения нагрузки на заданные УВ. Коммутация выполняется уставками в составе каждого элемента УВ на ОН. Для выполнения централизованного включения нагрузки в функции предусмотрен отдельный вход, что позволяет выполнять включение нагрузки одной командой.

\item
Функциональный блок может быть полностью выведен из работы оперативно путем активации сигнала «ОВ УОН».
\item
В таблице \ref{uvon:tbl1} приведены параметры, необходимые для настройки УОН.

\end{enumerate}

%%%%%%%%%%%%%%%%%%%%%%%%%%%%%%%%%%%%%%%%%%%%%%%%%%%%%%%%%  УВ ОН %%%%%%%%%%%%%%%%%%%%%%%%%%%%%%%%%%%%%%%%%%%%%%%%

\item Логика реализации управляющих воздействий на отключение нагрузки (УВ ОН1 ... УВ ОН10)\label{punktuvon}

\begin{enumerate}

\item
Ввод каждой функции УВ ОН в работу на этапе параметрирования устройства осуществляется своим программным переключателем SGF1 «Ввод\_функции». Информация о введенной функции отображается выходным сигналом для каждой функции «УВ ОН(m): Ввод», здесь и далее в этом подразделе: (m) -- номер функции, $\{m \in \mathbb{N} \mid 1 \leq m \leq 10\}$.

\item
Логика УВ ОН реализует 10 управляющих воздействий на отключение нагрузки. Реализация срабатывания ступени отключения выполнена в виде соответствующих входных дискретных сигналов команд отключения по 10 команд на каждое УВ ОН, объединенных общей логикой. Действие дискретных сигналов контролируются уставками SGF2...SGF11.

\item
Логика УВ ОН может быть выведена из работы оперативно путем активации сигнала вывода «ОВ УВ ОН(m)». Состояние логики характеризуется активным состоянием сигнала «УВ ОН(m): Оперативный вывод».

\item
При активном состоянии входного сигнала «УВ ОН(m) на сигн»  действие соответствующей логики УВ ОН переводится на сигнал.

\item
Алгоритм работы логики УВ ОН выполнен в соответствии с рисунком \ref{uvon:img1}. 

\vspace{3mm}
\begin{figure}[H]
\centering
\includegraphics[width=0.95\textwidth,height=0.95\textheight,keepaspectratio]{\fbpath/207. УОН/_latex/img/uvon.pdf}
\captionof{figure}{Функциональная схема логики УВ ОН}\label{uvon:img1}
\end{figure}

\end{enumerate}

%%%%%%%%%%%%%%%%%%%%%%%%%%%%%%%%%%%%%%%%%%%%%%%%%%%%%%%%%  УВ ЦВН %%%%%%%%%%%%%%%%%%%%%%%%%%%%%%%%%%%%%%%%%%%%%%%%

\item Логика реализации управляющих воздействий на централизованное включение нагрузки (УВ ЦВН1 ... УВ ЦВН10)\label{punktuvzvn}

\begin{enumerate}

\item
Ввод каждой функции УВ ЦВН в работу на этапе параметрирования устройства осуществляется своим программным переключателем SGF1 «Ввод\_функции». Информация о введенной функции отображается выходным сигналом для каждой функции «УВ ЦВН(m): Ввод».

\item
Логика УВ ОН реализует 10 управляющих воздействий на централизованное включение нагрузки.

\item
Реализация срабатывания ступени включения выполнена в виде одного входного дискретного сигнала «Команда ЦВН» для каждого УВ ЦВН, объединенных общей логикой. У каждого УВ ЦВН имеется своя выдержка времени на срабатывание функции, задаваемая уставкой Т1. Элементы УВ ЦВН для включения нагрузки позволяют задавать выдержки времени с целью ступенчатого включения отключенной нагрузкой.

\item
Логика УВ ЦВН может быть выведена из работы оперативно путем активации сигнала вывода «ОВ УВ ЦВН(m)». Состояние логики характеризуется активным состоянием сигнала «УВ ЦВН(m): Оперативный вывод».

\item
При активном состоянии входного сигнала «УВ ЦВН(m) на сигн»  действие соответствующей логики УВ ЦВН переводится на сигнал.

\item
Алгоритм работы логики УВ ЦВН выполнен в соответствии с рисунком \ref{uvzvn:img1}. 


\vspace{3mm}
\begin{figure}[H]
\centering
\includegraphics[width=0.95\textwidth,height=0.95\textheight,keepaspectratio]{\fbpath/207. УОН/_latex/img/uvzvn.pdf}
\captionof{figure}{Функциональная схема логики УВ ЦВН}\label{uvzvn:img1}
\end{figure}

\end{enumerate}

%%%%%%%%%%%%%%%%%%%%%%%%%%%%%%%%%%%%%%%%%%%%%%%%%%%%%%%%% ОБЩИЕ ЦЕПИ УОН %%%%%%%%%%%%%%%%%%%%%%%%%%%%%%%%%%%%%%%%%%%%%%%%

\item Сборка сигналов УОН \label{punktssuon}

\begin{enumerate}

\item
При объединении управляющих воздействий УОН формируются следующие выходные сигналы:
\begin{itemize}
    \item	«УОНобщ: Срабатывание ОН» –- формируется при срабатывании любого УВ ОН, в том числе переведенного в работу только на сигнал;
    \item	«УОНобщ: Срабатывание ЦВН» –- формируется при срабатывании любого УВ ЦВН, в том числе переведенного в работу только на сигнал;
    \item	«УОНобщ: Срабатывание» -- формируется при наличии одного из сигналов «УОНобщ: Срабатывание ОН» или «УОНобщ: Срабатывание ЦВН».
\end{itemize}

\item
Алгоритм работы функции сборки сигналов УОН выполнен в соответствии с рисунком \ref{ssuon:img1}.

\vspace{3mm}
\begin{figure}[H]
\centering
\includegraphics[width=0.7\textwidth,height=0.7\textheight,keepaspectratio]{\fbpath/207. УОН/_latex/img/uoncommon.pdf}
\captionof{figure}{Функциональная схема сборки сигналов УОН}\label{ssuon:img1}
\end{figure}

\end{enumerate}

\end{enumerate}
